% TOP_PROBLEM: {"id": 3132, "title": "Minimum number of arc-disjoint transitive subtournaments of order 3 in a tournament", "category_id": 3, "category": {"id": 3, "name": "graph_theory", "display_name": "Graph Theory", "description": "Problems involving graphs, networks, and their properties.", "slug": "graph-theory", "order_index": 3, "created_at": "2026-07-31T15:26:25.671Z"}, "status": "open", "problem_number": "OPG-49795", "set_id": 12, "statusQualification": "The auxiliary extremal function is the minimum, over tournaments, of the maximum packing size; the source discussion omits this inner maximum."}
% FIRST_POSTED: 2026-10-01
% ATTEMPT_STATUS: unresolved
% UnsolvedMath ID 3132; source catalogue code OPG-49795
% !TeX program = lualatex
% Research attempt by GPT-6 Astra Ultra; no claim of a new complete solution.
\documentclass[11pt,a4paper]{article}
\usepackage[margin=25mm]{geometry}
\usepackage{fontspec}
\setmainfont{lmroman10-regular.otf}[BoldFont=lmroman10-bold.otf,
 ItalicFont=lmroman10-italic.otf,BoldItalicFont=lmroman10-bolditalic.otf]
\usepackage{amsmath,amssymb,mathtools}
\usepackage{unicode-math}
\setmathfont{latinmodern-math.otf}
\AtBeginDocument{\let\setminus\smallsetminus}
\usepackage{xurl,hyperref}
\hypersetup{colorlinks=true,urlcolor=blue}
\setcounter{secnumdepth}{0}
\setlength{\parindent}{0pt}
\setlength{\parskip}{5pt}
\setlength{\emergencystretch}{3em}
\begin{document}
\section*{Minimum number of arc-disjoint transitive subtournaments of order 3 in a tournament}\label{3132}
{\small UnsolvedMath ID 3132 · OPG-49795 · Graph Theory · OpenGarden}
\subsection{Definitions and mathematical statement}
A tournament of order $n\geq1$ is a directed graph with $n$ vertices,
no loops, and exactly one of the arcs $u\to v$ and $v\to u$ for each
pair of distinct vertices. A transitive triple is a three-vertex
subtournament whose vertices can be ordered $a,b,c$ with arcs
$a\to b$, $a\to c$, and $b\to c$. Thus its three arcs do not form
a directed cycle.

A packing of transitive triples is a collection of such triples whose
arc sets are pairwise disjoint. Different triples may share a vertex, but no pair can share
two vertices, but cannot reuse any directed edge. Write $\nu_T(T)$
for the largest number of triples in a packing in $T$ and define
\[
 f(n)=\min\{\nu_T(T):T\text{ is a tournament on }n\text{ vertices}\}.
\]
The maximum inside this minimum is essential: an arbitrary packing
could always be empty.

The conjectured guarantee is that every tournament $T$ of order $n$
satisfies
\[
 \nu_T(T)\geq
 \left\lceil\frac{n(n-1)}6-\frac n3\right\rceil
 =\left\lceil\frac{n(n-3)}6\right\rceil.
\]
The ceiling denotes the least integer greater than or equal to its
argument. The assertion asks for one simultaneous packing of this
size; counting all transitive triples without controlling their
overlap is insufficient. Since each selected triple uses three arcs,
the scale of the proposed bound is one third of all arcs in the
tournament, with a linear correction.
\subsection{Short English statement}
Can every tournament pack at least the stated number of transitive triples without reusing an arc?
\subsection{Research attempt}
\paragraph{Scope and process.}
This is a GPT-6 Astra Ultra research attempt dated 1 October 2026, using
primary-source browsing and elementary mathematical checks. The canonical
definition above is retained. Results below are restricted results or
reductions, not a claim of a new solution to the entire catalogue question.
No formal proof assistant or external mathematical referee checked this text.


\paragraph{Literature and approach.}
The original source [S2, R1] attributes the conjecture to Yuster and
describes a balanced cyclic three-part construction showing the proposed
bound is sharp. We verify that obstruction, then give a self-contained
packing guarantee from a maximal packing. The guarantee is weaker than
the conjecture and than sophisticated published packing bounds; its
purpose is to isolate the loss in a simple greedy attack.

\paragraph{A sharp obstruction to any stronger proposed bound.}
Partition the vertices into $V_1,V_2,V_3$ as equally as possible and
orient all cross-arcs cyclically, $V_1\to V_2\to V_3\to V_1$.
Orient within the parts arbitrarily. A triple with one vertex in
each part is cyclic, so every transitive triple uses an arc internal
to some $V_i$. Arc-disjoint triples use distinct such internal arcs.
Consequently every packing has size at most
\[
 \sum_{i=1}^3\binom{|V_i|}{2}
 =\left\lceil\frac{n(n-3)}6\right\rceil
 \qquad(n\ge2).
\]
For the identity, put $n=3q+r$ with $0\le r<3$ and use $r$ parts
of size $q+1$ and $3-r$ of size $q$. This proves the upper obstruction,
not existence of a packing reaching it in every orientation.
For $n=1$ both the actual packing number and the proposed ceiling
are zero, separately.

\paragraph{What a maximal packing leaves behind.}
Start with any packing to which no additional arc-disjoint transitive
triple can be added. Delete its $3k$ arcs and ignore the directions
of the remaining arcs, obtaining a simple graph $R$. This graph
contains no $K_4$. Indeed, if four vertices retained all six arcs,
their induced tournament would contain a transitive triple: at any
one vertex at least two of the other three are both out-neighbours
or both in-neighbours, and those three vertices form such a triple
regardless of the direction of their remaining arc. It could have
been added to the packing, contradicting maximality.

\paragraph{An elementary bound for the residual graph.}
We prove that a $K_4$-free graph on $n$ vertices has at most $n^2/3$
edges. If it has no triangle, every edge $uv$ satisfies
$d(u)+d(v)\le n$, because the two neighbourhoods are disjoint.
Summing over edges gives $\sum_vd(v)^2\le ne$; Cauchy--Schwarz
gives $4e^2/n\le\sum_vd(v)^2$, and hence $e\le n^2/4$
(the case $e=0$ is immediate). Otherwise choose a triangle. Every
outside vertex has at most two neighbours in it, so deleting the
triangle removes at most $3+2(n-3)=2n-3$ edges. Induction on $n$
now gives
\[
 e\le(n-3)^2/3+(2n-3)=n^2/3.
\]
The base cases $n\le3$ satisfy the bound directly.

Apply this to $R$. Since $|E(R)|=\binom n2-3k$,
\[
 k\ge\left\lceil\frac{n(n-3)}{18}\right\rceil.
\]
Any inclusion-maximal packing meets this weaker guarantee, so the
maximum packing does too. The conjecture would require roughly
three times as many triples; the residual extremal bound alone
therefore leaves a substantial quadratic gap.

\paragraph{Finite verification and what it establishes.}
The exact script
\texttt{scripts/check\_attempt\_batch02\_algebra\_2026\_10\_01.py}
enumerates all orientations for $n\le5$. For each it lists the
transitive three-vertex sets and computes the maximum collection of
pairwise disjoint arc masks by exhaustive recursion. The resulting
minimum packing numbers for $n=1,2,3,4,5$ are $0,0,0,1,2$,
matching the proposed ceilings. This checks the small instances
and the maximum/minimum quantifiers, not the asymptotic conjecture.

\paragraph{Verification and first remaining gap.}
Maximal does not mean maximum; only the absence of an additional
disjoint triple was used in the residual argument. Removing three
arcs per chosen triple is legitimate because their arc sets are
disjoint, even though they can share vertices. The obstruction
construction counts a necessary internal arc, not an arbitrary
three-vertex set. To improve the lower bound to the proposed one,
one needs an exchange or absorption argument that replaces selected
triples and controls the residual orientation. Merely counting all
available triples or declaring the residual $K_4$-free cannot do so.

\paragraph{Final status.}
\textbf{UNRESOLVED.} The sharp upper obstruction, an explicit weaker
universal packing bound, and the finite cases through order five
are verified; no all-order matching lower bound is proved.

\subsection{Sources}
\begin{itemize}
\item[S1] ULAM.AI and the UnsolvedMath Contributors, supplied problems.json, record 3132 (OPG-49795), OpenGarden collection. Catalogue identity and source transcription; CC BY 4.0.
\url{https://huggingface.co/datasets/ulamai/UnsolvedMath}.
\item[S2] Open Problem Garden, Minimum number of arc-disjoint transitive subtournaments of order 3 in a tournament. Original problem and discussion; the mathematical formulation below is expanded from the supplied source record.
\url{http://www.openproblemgarden.org/op/minimum_number_of_transitive_subtournaments_of_order_3_in_a_tournament}.

\item[R1] Raphael Yuster, original conjecture as posted by Fr\'ed\'eric
Havet, \emph{Minimum number of arc-disjoint transitive subtournaments
of order 3 in a tournament}, original problem page and bibliography,
checked 1 October 2026.
\url{https://www.openproblemgarden.org/op/minimum_number_of_transitive_subtournaments_of_order_3_in_a_tournament}.

\end{itemize}
\end{document}
